\documentclass[10pt,a4paper]{amsart}
\usepackage[margin=19mm]{geometry}
\usepackage{amsmath,amssymb,mathtools}
\usepackage[scr=rsfs,cal=euler]{mathalfa}
\usepackage{xfrac}
\usepackage[unicode,hidelinks,bookmarks=false]{hyperref}
\newcommand{\ie}{i.e.\ }
\newcommand{\cf}{cf.\ }
\newcommand{\resp}{respectively\ }
\newcommand{\Subsection}{\subsection}
\providecommand{\nobreakdash}{\nobreak\hbox{-}}
\setlength{\emergencystretch}{2em}
\multlinegap0pt
\usepackage{bm}
\usepackage{bbm}
\usepackage{dsfont}
\usepackage{xspace}
\usepackage{cancel}
\usepackage{mathtools}
\usepackage{amsthm}
\makeatletter
\@ifundefined{lemma}{\newtheorem{lemma}{Lemma}}{}
\@ifundefined{proposition}{\newtheorem{proposition}[lemma]{Proposition}}{}
\@ifundefined{remark}{\newtheorem{remark}[lemma]{Remark}}{}
\makeatother
\newcommand{\R}{{\mathbb R}}
\newcommand{\proofend}{$\Box$\bigskip}
\title[Outer symplectic billiards]{Outer symplectic billiards (Proposition 3)}
\author{Peter Albers, Ana Chavez Caliz, Serge Tabachnikov}
\date{Source: arXiv:2409.07990v2 (CC BY 4.0)}
\begin{document}
\maketitle

\noindent\emph{Source and licence.} Outer symplectic billiards. Version arXiv:2409.07990v2 (CC BY 4.0), distributed under the Creative Commons Attribution 4.0 International licence, as recorded in the arXiv licence metadata for this exact version and shown on the abstract page https://arxiv.org/abs/2409.07990v2. Full rights notice: licenses/arxiv-2409.07990-CC-BY-4.0.txt.\par\smallskip

\noindent\textit{Editorial scope.} Counted unit: the Remark printed as Proposition 3 of the article (source0.tex lines 1697--1703, source label \texttt{None}), delivered together with the article's own complete original proof of it (source0.tex lines 1705--1780, one \texttt{proof} environment of 76 lines). No mathematics is written, repaired or re-derived: statement and proof are the article's own text line for line. Required context retained uncounted: lem:description\_of\_wall\_Lagrangian\_case (lines 1267--1280); prop:abcd\_example\_and\_wall (lines 1550--1566). Every definition, display and label of those lines is retained unchanged. Omitted from this package and not redistributed: 1--1266, 1281--1549, 1567--1696, 1704--1704, 1781--2370. Nothing inside a retained excerpt is omitted or altered: every retained body is delivered exactly as the article's lines are written, so no omission receipt of any kind accompanies this package. Citations: the retained excerpts carry no citation command at all, so no bibliography entry of the article is transcribed and no reference is redistributed. Reference adaptation: none was needed and none was made; every reference inside the delivered unit resolves inside it, so no unresolved reference or citation is produced. Printed slips kept literal: no printed word, symbol, hypothesis or number of the counted statement or of its proof was altered. Layout and class: the article's own class is \texttt{article} with packages including epsfig, graphics, amsmath, amsfonts, latexsym, amssymb, graphicx, url, epstopdf, hyperref, xcolor, comment, mathrsfs; the wrapper is the lane's pinned \texttt{amsart} editorial wrapper at 10pt on A4, which restates the theorem-like environments the retained text needs and adds the article's own macro definitions that the retained text needs (\texttt{\textbackslash R}, \texttt{\textbackslash proofend}), with extra wrapper packages (bm, bbm, dsfont, xspace, cancel, mathtools). The delivered numbering is the unit's own. Assets: no retained span contains a figure environment, a graphics-inclusion command, a PDF-page inclusion, a table or a TikZ picture; no image file is copied into this package, the package declares no asset dependency and no image file is redistributed with it. Licence: version arXiv:2409.07990v2 of this article is distributed under the Creative Commons Attribution 4.0 International licence, as recorded in the arXiv licence metadata for this exact version and shown on the abstract page https://arxiv.org/abs/2409.07990v2. A full rights notice is delivered at licenses/arxiv-2409.07990-CC-BY-4.0.txt. Characters: every retained excerpt is pure ASCII, so the delivered engine, XeLaTeX with its default Latin Modern text font, reports no missing character.
\begin{lemma}\label{lem:description_of_wall_Lagrangian_case}
The set $\Delta$ of singular points of $\Psi:\R^n\times\R^n\to \R^{2n}$ is 
\begin{equation}\nonumber
\begin{aligned}
\Delta=\big\{(q,w)\in\R^n\times\R^n\mid\R^n\ni\zeta\mapsto \nabla^3F(q)[\zeta,w]\in\R^n\text{ is singular}\,\big\}.
\end{aligned}
\end{equation}
%
Therefore, a point $\Psi(q,w)=(q+w,\nabla F(q)+\nabla^2F(q)w)\in\R^{2n}$ is an element of $\Sigma$ if and only if the linear map
\begin{equation}\nonumber
\R^n\ni\zeta\mapsto \nabla^3F(q)[\zeta,w]\in\R^n
\end{equation}
does not have full rank. 
\end{lemma}

\begin{proposition}\label{prop:abcd_example_and_wall} 
The following statements are equivalent:
\begin{enumerate}   \renewcommand{\theenumi}{\roman{enumi}}\renewcommand{\labelenumi}{\rm (\theenumi)} %% \itemsep=1.5ex
\item The coefficients $a,b,c,d$ of $F$ satisfy the inequality 
\begin{equation}\label{eq:no_wall_condition_cubic_examples}
18abcd+b^2c^2>4ac^3+4b^3d+27a^2d^2.
\end{equation}
\item $\Delta=L$, where $L\subset TL$ denotes the zero-section, and then, in particular, $\Sigma=L\subset\R^4$.
%\item The multiplicity of the map $\Psi:TL\to\R^{2n}$ on $\R^{2n}\setminus L$ is 2 everywhere. That is, $\Psi:TL\setminus L\to\R^{2n}\setminus L$ is a local diffeomorphism and each point in $\R^{2n}\setminus L$ has precisely two distinct preimages.
\item Every regular value of the map $\Psi:TL\to\R^4$ has multiplicity 2.
\end{enumerate}
The opposite inequality 
\begin{equation}\label{eq:no_wall_condition_cubic_examples_other_way_round}
18abcd+b^2c^2<4ac^3+4b^3d+27a^2d^2,
\end{equation}
is equivalent to the multiplicity of $\Psi$ being $0$ and $4$. In this case, $\Delta\neq L$.
\end{proposition}

\begin{remark}
{\rm
Certainly, the statement that $L$ contains a line is independent of a generating function. We can rephrase this as the existence of a vector $W\in\R^{2n}$ such that $X+\R W\subset L$, and of course, then $W\in T_XL$.

If any of the equivalent assertions in 1.~is true, the equality $\Delta=TL$ follows from combining Lemma \ref{lem:description_of_wall_Lagrangian_case} with $\nabla^3F(q)[w,\zeta]=\nabla^2F(w)\zeta$. In particular, then $\Sigma=\Psi(\Delta)=\Psi(TL)$, and the outer symplectic billiard correspondence is not well-defined anywhere.
} 
\end{remark}

\begin{proof}
It is enough to prove the equivalence between 1.(b) and 1.(c) since the latter condition is actually independent of the point $q$. That is, then 1.(b) holds for all $q$ with the same vector $w$ which is 1.(a).

Let us assume that $L$ contains a line, i.e., we can find $q,w\in\R^n$ with $w\neq0$ such that
\begin{equation}\nonumber
\Psi(q,tw)=\big(q+tw,\nabla F(q)+t\nabla^2F(q)w\big)\in L\quad\forall t\in\R.
\end{equation}
In particular, we find $Q(t)\in\R^n$ such that
\begin{equation}\nonumber
\big(q+tw,\nabla F(q)+t\nabla^2F(q)w\big)=\big(Q(t),\nabla F(Q(t))\big).
\end{equation}
We immediately obtain $Q(t)=q+tw$. The Taylor expansion 
\begin{equation}\nonumber
\nabla F(q+tw)=\nabla F(q) +t \nabla^2F(q)w+\tfrac12t^2\nabla^3F(q)[w,w]
\end{equation}
is exact since $F$ is cubic, and it follows that
\begin{equation}\nonumber
\nabla^3F(q)[w,w]=0.
\end{equation}
Using again the equality $\nabla^3F(q)[w,w]=\nabla^2F(w)w$, we see that if $L$ contains a line then $\nabla^2F(w)w=0$ for some $w\neq0$. The above computation also shows that $\nabla^2F(w)w=0$ implies that $\Psi(q,tw)\in L$ for all $t\in\R$. This proves point 1.~in the proposition.

To show 2.~we recall from the proof of Proposition \ref{prop:abcd_example_and_wall} that
\begin{equation}\nonumber
\nabla^2F(w)w=\begin{pmatrix}
6aw_1+2bw_2 & 2bw_1+2cw_2\\
2bw_1+2cw_2 & 2cw_1+6dw_2
\end{pmatrix}
\begin{pmatrix}
w_1\\
w_2
\end{pmatrix}
=0
\end{equation}
are the two real equations
%\begin{equation}\nonumber
%\left\{\;
%\begin{aligned}
%(3aw_1+bw_2)w_1+(bw_1+cw_2)w_2&=0\\
%(bw_1+cw_2)w_1+(cw_1+3dw_2)w_2&=0
%\end{aligned}
%\right.
%\end{equation}
\begin{equation}\nonumber
\left\{\;
\begin{aligned}
3aw_1^2+2bw_1w_2+cw_2^2&=0\\
bw_1^2+2cw_1w_2+3dw_2^2&=0.
\end{aligned}
\right.
\end{equation}
Again $w_1=w_2=0$ is a solution and if $w_2\neq0$ this systems reduces to the quadratic equations in $z=\frac{w_1}{w_2}$ given by
\begin{equation}\nonumber
\left\{\;
\begin{aligned}
3az^2+2bz+c&=0\\
bz^2+2cz+3d&=0.
\end{aligned}
\right.
\end{equation}
A classical fact is that two complex polynomials in one variable have a common root if and only if their resultant is zero. The resultant of the two equations from above is by definition
\begin{equation}\nonumber
\det
\begin{pmatrix}
3a & 0 & b & 0 \\
2b & 3a & 2c & b \\
c & 2b & 3d & 2c \\
0 & c  & 0 & 3d 
\end{pmatrix}=3\big(18abcd+b^2c^2-(4ac^3+4b^3d+27a^2d^2)\big).
\end{equation}
Since we are working over reals, for the quadratic equations to have solutions we also need their discriminants to be positive:
$$
b^2 - 3ac > 0,\ c^2 - 3bd > 0. 
$$ 
We conclude that $\nabla^2F(w)w=0$ for some $w\neq0$ is equivalent to the claimed equation $18abcd+b^2c^2=4ac^3+4b^3d+27a^2d^2$ and the inequalities $b^2>3ac, c^2>3bd$.
\proofend 
\end{proof}

\end{document}
